\subsection{Discussion}\label{sec:discussion}

The findings provide an economic interpretation of how firms' economic
footprints can evolve in response to changing market opportunities.
Distribution-valued representations allow researchers to compare firms
through the information associated with their activities, but changes in
those distributions do not, by themselves, explain the economic decisions
underlying repositioning.
Our model connects these movements to strategic competition.
When firms adjust capacity across overlapping activities, their decisions
affect competitors' margins, while adjustment costs make the response to new
opportunities gradual.
In the restricted two-opportunity economy, initial operating exposures
determine both the ordering of return responses and the direction of
movement relative to competitors' initial news distributions.
This establishes a mechanism through which competitive adjustment can
generate systematic changes in firms' reported economic footprints, extending
distribution-based comparisons beyond static measures of similarity.

The distinction between financial valuation and operating adjustment is
central to this mechanism.
Investors can capitalize the anticipated cash-flow consequences of an
exogenous free cash flow shock when information arrives, whereas firms adjust their
operating allocations over subsequent periods.
Return responses and changes in reported economic footprints consequently need not
coincide in time.
The article-timing weights in our estimator connect reference firms' daily
excess returns to the subsequent appearance of information about the
adjusting firm.
This construction captures a directional relationship between weighted
excess returns and changes in the firm's reported economic footprint, rather
than identifying a structural adjustment speed.
Its economic interpretation depends jointly on initial exposures, competitive
interactions, adjustment costs, and the timing of information.

The historical results reveal a positive but economically small association.
The estimated direction is consistent across four measurement specifications,
and the primary normalized association is \ppnum[4]{ir-normalized-alignment},
corresponding to approximately
\num[round-mode=places,round-precision=2]{\fpeval{100*\ppvalue{ir-d2560-expanded-within-fit}}}\%
descriptive within-sample fit.
Thus, the predicted ordering appears in the data but accounts for only a
limited component of observed information movements.
This modest magnitude is compatible with gradual operating adjustment,
heterogeneous responses to economic opportunities, and imperfect measurement
of firms' activities.
The simulations clarify the last possibility: at the benchmark's small
adjustment scale, finite article samples frequently fail to resolve a
positive population relationship, even when coverage increases substantially.
They do not, however, establish which mechanism explains the small historical
association.
The evidence supports a directional correspondence without demonstrating that
the underlying operating response is either strong or uniform across firms.

The reconstruction results sharpen this interpretation.
In the primary specification, the restricted movement profile improves
held-out reconstruction of annual news changes by only
\ppnum[2]{ir-observation-primary-heldout-pct}\% relative to a zero-change
benchmark, while two alternative specifications perform worse than that benchmark.
A positive directional association can therefore coexist with weak
reconstruction of the broader evolution of firms' reported economic footprints.
These findings answer different questions.
The directional estimate measures whether relative movements toward reference
firms correspond to their relative weighted excess returns, whereas the reconstruction
exercise evaluates how well the model's restricted geometric pattern
describes observed changes in news distributions.
The theory predicts the former under specified conditions, not a complete
account of the latter.
Nevertheless, the weak reconstruction indicates that the historical movements
contain substantial variation beyond the local theoretical approximation.
The empirical evidence is consequently more informative about relative
directional ordering than about the full geometry or magnitude of firms' repositioning.

An additional implication concerns the production of financial information itself.
Corporate news reflects both underlying economic activity and the attention
directed toward different activities.
Under the observation model, changes in reporting intensity alter a firm's
information distribution even when its operating allocation remains unchanged.
The reporting-only simulation demonstrates that this channel can generate
positive directional alignment at fixed capacities.
Operating adjustment and reporting responses can therefore produce similar
observed movements, potentially reinforcing or offsetting one another.
This distinction is not simply a source of statistical noise: it affects the
economic meaning of the distances used to represent firms.
The historical estimator measures changes in reported economic footprints
associated with relative weighted excess returns, but does not separately identify the
contributions of operating decisions and reporting intensity.
Interpreting information-space movements as changes in economic activity
requires assumptions about how that activity becomes observable.

The stability of the estimated sign provides useful evidence about the
historical association, although it does not establish its generality.
The primary coefficient remains positive across annual transitions and the
reported deletion exercises.
At the same time, the 2020 transition contributes
\ppnum[1]{ir-primary-2020-weight-pct}\% of the pooled centered financial
second moment, and the largest reference-firm contribution accounts for
\ppnum[1]{ir-reference-share-pct}\%.
The financial variation informing the estimate is therefore unevenly
distributed across the sample.
Moreover, the theoretical prediction arises under restricted opportunity and
observation conditions that need not characterize the wider corporate environment.
Stability of the estimated direction should accordingly be distinguished from
evidence that the same economic mechanism operates with comparable strength
across firms, activities, and market conditions.

Taken together, these results extend the economic interpretation of
distribution-valued firms by linking changes in reported economic footprints to
competitive adjustment and financial valuation.
Building on the multi-contest allocation framework of
\citet{xu_zenou_zhou_conflict_2022}, the analysis establishes how inherited
capacities, strategic responses, and anticipated cash flows can jointly
produce a testable directional relationship in information space.
The historical application identifies a corresponding association, while the
simulations and reconstruction exercises establish important boundaries on
its interpretation.
More direct tests would combine dated information about investment,
production, or operating capacity with financial returns and corporate news,
allowing changes in economic activity to be distinguished from changes in reporting.
Such evidence could also examine whether the predicted ordering extends
beyond the restricted setting considered here.
The broader contribution is an economic foundation for studying firms'
changing economic footprints through financial news.
